\section*{A paid rectangle interface for ratio-critical Boolean sources}

All Fourier entropies are in bits.  Let
$f:\{-1,1\}^{A\sqcup B}\to\{-1,1\}$ be Boolean, with
$|A|=a$, $|A|+|B|=N$.  Define the multiplicative rectangle defect
\[
 \rho_A(f)=\Pr_{x,x'\in\{-1,1\}^A,\ y,y'\in\{-1,1\}^B}
 \bigl[f(x,y)f(x',y)f(x,y')f(x',y')=-1\bigr].       \tag{1}
\]
This is a truth-table statistic fixed before any entropy comparison.  It is
zero exactly when $f$ is a disjoint Boolean product across the cut.

\subsection*{The anchoring identity}

For an anchor $(x_0,y_0)$ put
\[
 u_{y_0}(x)=f(x,y_0),\qquad
 v_{x_0,y_0}(y)=f(x_0,y)f(x_0,y_0),
\]
and $T_{x_0,y_0}(x,y)=u_{y_0}(x)v_{x_0,y_0}(y)$.  Then
\[
 \mathbb E_{x_0,y_0}\Pr_{x,y}[f(x,y)\ne T_{x_0,y_0}(x,y)]
 =\rho_A(f).                                         \tag{2}
\]
Indeed, the mismatch indicator is one precisely when the product of the four
corners $(x,y),(x,y_0),(x_0,y),(x_0,y_0)$ is $-1$.  The four coordinates are
independent and uniform after averaging the query and anchor.  Consequently
there are Boolean $u,v$ such that
\[
 \varepsilon:=\Pr[f\ne uv]\le\rho_A(f).             \tag{3}
\]
This proves a Boolean realizability statement, not merely a rank-one
approximation of an arbitrary spectral law.

\subsection*{A dimension-explicit component excess}

We first record a harmless universal estimate.  If $g$ is nonconstant on
$N$ bits and $R(g)=H(g)/I(g)$, then
\[
 R(g)\le3N.                                           \tag{4}
\]
To prove this, let $S$ have law $\widehat g(S)^2$ and
$q=\Pr[S\ne\varnothing]=\operatorname{Var}(g)$.  Then
$I(g)=\mathbb E|S|\ge q$ and
\[
 H(g)\le h_2(q)+Nq.
\]
Writing the minority truth-table density as $\delta$, nonconstancy gives
$\delta\ge2^{-N}$ and $q=4\delta(1-\delta)\ge2^{1-N}$.  The elementary bound
$h_2(q)\le q\log_2(1/q)+q/\ln2$ therefore gives
$h_2(q)\le2Nq$, which proves (4).

For $0\le r\le1/16$ define the interface excess
\[
 \mathcal E_{N,a}(r)
 =2h_2(2\sqrt r)+(2N-a)2\sqrt r+6aNr.                \tag{5}
\]
It depends only on the physical cut, its rectangle defect, and the ambient
dimension.  It is not defined as a missing entropy gap.

\begin{theorem}[controlled component excess]
Suppose $f$ is descendant-critical: every nonconstant physical restriction
$g$ satisfies $R(g)\le R:=R(f)$.  If $\rho=\rho_A(f)\le1/16$, then there are
Boolean $u$ on $A$ and $v$ on $B$, with
$\varepsilon=\Pr[f\ne uv]\le\rho$, for which
\[
 \boxed{\quad R I(u)-H(u)\le\mathcal E_{N,a}(\rho).\quad}    \tag{6}
\]
In particular, when $u$ is nonconstant,
\[
 (R-R(u))I(u)\le\mathcal E_{N,a}(\rho).              \tag{7}
\]
The symmetric assertion holds for $v$, replacing $a$ by $N-a$.
\end{theorem}

\begin{proof}
Choose $u,v$ from (3).  The A053 critical near-tensor theorem, applied to the
same physical cut, gives
\[
 R\bigl(I(u)-2a\varepsilon\bigr)
 \le H(u)+2h_2(2\sqrt\varepsilon)
             +(2N-a)2\sqrt\varepsilon.               \tag{8}
\]
Because $\varepsilon\le\rho\le1/16$, the displayed continuity terms are
nondecreasing when $\varepsilon$ is replaced by $\rho$.  Move $H(u)$ to the
left, use $R\le3N$, and bound
$2Ra\varepsilon\le6aN\rho$.  This is (6).  The proof for $v$ is identical.
\end{proof}

\subsection*{An explicit influence-paid class}

The preceding excess is not automatically $O(I)$ for arbitrary cuts.  The
following sufficient admission criterion states exactly when it is paid.

\begin{corollary}[paid low-interface cut]
Let $J=I_A(f)>0$, $0<\kappa\le1$, and put $t=2\sqrt{\rho_A(f)}$.  Suppose
\[
 t\le\min\left\{
 \frac{\kappa^2J^2}{144},
 \frac{\kappa J}{6N},
 \sqrt{\frac{2\kappa J}{9aN}}
 \right\}.                                           \tag{9}
\]
Then the factors in the theorem may be chosen with
\[
 I(u)\ge\left(1-\frac{\kappa}{9N}\right)J\ge\frac89J, \tag{10}
\]
and the interface is paid by actual physical block influence:
\[
 R I(u)-H(u)\le\kappa J.                              \tag{11}
\]
Consequently
\[
 R(u)\ge R-\frac98\kappa.                            \tag{12}
\]
\end{corollary}

\begin{proof}
Hamming perturbation of a product changes each block influence by at most
$2a\varepsilon$, so
$I(u)\ge J-2a\varepsilon$.  From the third condition in (9),
$2a\varepsilon=a t^2/2\le\kappa J/(9N)$, proving (10).
The elementary inequality $h_2(t)\le2\sqrt t$ and the three conditions in
(9) respectively give
\[
 2h_2(t)\le\frac{\kappa J}{3},\qquad
 (2N-a)t\le\frac{\kappa J}{3},\qquad
 6aN\rho=\frac32aNt^2\le\frac{\kappa J}{3}.
\]
Thus $\mathcal E_{N,a}(\rho)\le\kappa J$, and (11)--(12) follow.
\end{proof}

For $\rho_A(f)=0$, the result specializes to exact tensor purity: $f=uv$ and
both nonconstant factors have ratio exactly $R$.  For positive defect, it is a
quantitative alternative: either the cut violates the explicit small-interface
condition (9), or it exposes a Boolean factor whose ratio is within
$9\kappa/8$ of the critical source.  Universal admission of some cut satisfying
(9) remains open.
